\capredactadotrue
\section{Plan Financiero-Económico}

Se determina el presupuesto financiero de \nombreproyecto{} para un horizonte temporal de tres años (36 meses). Este horizonte exige un modelo de adopción acelerada y una política de austeridad en los primeros ejercicios, para que el proyecto alcance su rentabilidad, medida a través del valor actual neto (VAN), la tasa interna de retorno (TIR) y el período de recupero, dentro del plazo previsto. Todos los importes se expresan en dólares estadounidenses (USD), aunque los cobros puedan emitirse en pesos argentinos. Los flujos se presentan antes de impuestos (véanse las secciones sobre impuestos).

\subsection{Modelo de ingresos}
La plataforma opera bajo un modelo \textit{freemium} B2C complementado con un \textit{marketplace} C2C. Por el acotamiento del horizonte a tres años y la necesidad de cubrir los costos operativos de un equipo de hasta diez integrantes, los ingresos requieren un escalado acelerado bajo un esquema de precios híbrido.

Se proyectan ingresos de USD 35.000 en el año 1, USD 95.000 en el año 2 y USD 195.000 en el año 3. Para lograrlo sin perder la accesibilidad regional, la suscripción anual premium se escalona en USD 35 (año 1), USD 50 (año 2) y USD 70 (año 3), lo que exige captar aproximadamente 986, 1.889 y 2.775 suscriptores, respectivamente.

La apertura teórica de las ventas establece que el 100\,\% ingresa al contado a través de pasarelas de pago digitales.

\subsection{Modelo de egresos}
Los egresos se desglosan según la evolución operativa.

\subsubsection{Costos variables (33\,\% de los ingresos)}
Están directamente ligados a los ingresos y comprenden el presupuesto de marketing digital (15\,\%), las comisiones de las plataformas de descarga y de pago (15\,\%) y las comisiones bancarias (3\,\%). Ascienden a USD 11.550 en el año 1, USD 31.350 en el año 2 y USD 64.350 en el año 3.

\subsubsection{Costos fijos operativos}
Comprenden el capital de trabajo anual presupuestado en USD 14.400, constante, y los costos de licencias de tecnologías de la información, de mantenimiento de la infraestructura propia (energía, conectividad, dominio y respaldos), de un eventual escalado a servicios en la nube y de seguros. Estos últimos aumentan de USD 1.000 (año 1) a USD 1.500 (año 2) y USD 2.500 (año 3), a medida que crece la carga sobre los servidores. Las remuneraciones y cargas sociales se detallan en la sección correspondiente. No se contempla la contratación de deuda bancaria.

\subsection{Modelo de inversión}
La inversión inicial necesaria antes de comenzar a operar es de USD 27.610 y se compone de:
\begin{itemize}
  \item \textbf{Activos fijos (USD 3.360):} compra de infraestructura y equipamiento informático usado (servidores y \textit{notebooks}), para minimizar el impacto inicial.
  \item \textbf{Activos diferidos (USD 1.500):} registros legales y publicidad de lanzamiento.
  \item \textbf{Capital de trabajo (USD 22.750):} fondo de maniobra para afrontar los primeros meses de operación hasta consolidar la base de suscriptores activos.
\end{itemize}
Como el lanzamiento comercial se produce en el mes 6, los primeros cinco meses transcurren sin ingresos y su costo operativo ronda los USD 7.500. El resto del capital de trabajo actúa como reserva de contingencia ante desvíos del escenario base.

\subsection{Amortizaciones y depreciaciones}
Los activos fijos y diferidos se deprecian y amortizan por el método de línea recta en el horizonte de 36 meses: los activos fijos (USD 3.360) a razón de USD 1.120 por año y los activos diferidos (USD 1.500) a razón de USD 500 por año, es decir, USD 1.620 anuales en total. Se trata de cargos contables sin salida de fondos, que reducen la base imponible del impuesto a las ganancias.

\subsection{Impuesto a las ganancias}
El impuesto a las ganancias se estima con una alícuota general del 25\,\%. Los flujos de este capítulo se presentan antes de impuestos; el efecto de este tributo sobre la caja se determinará junto con el equipo contable designado, según el resultado imponible de cada ejercicio y las amortizaciones descritas.

\subsection{Impuesto a los ingresos brutos}
El impuesto a los ingresos brutos se liquida mensualmente, aplicando la alícuota vigente en la jurisdicción provincial de residencia del proyecto (Santa Fe) sobre las ventas netas. Como la alícuota varía entre provincias y según la actividad, su determinación y su tratamiento se encomiendan al contador o equipo contable designado, y no se incluyen en los flujos de este capítulo.

\subsection{Remuneraciones y cargas sociales}
La plantilla de recursos humanos explica el principal incremento de los egresos y se divide en las tres etapas del plan.
\begin{itemize}
  \item \textbf{Lanzamiento (año 1).} La estructura mínima, de dos fundadores, genera un costo nominal en sueldos de USD 12.000. Mediante la política de reinversión del sueldo emprendedor, el 100\,\% de ese importe se inyecta nuevamente en el proyecto, de modo que no implica salida de fondos. El único egreso real son las cargas sociales de directores (USD 2.500).
  \item \textbf{Consolidación (año 2).} Con el equipo ampliado a cinco personas, la nómina de los tres nuevos profesionales alcanza los USD 24.000 en sueldos brutos y USD 6.000 en cargas sociales.
  \item \textbf{Expansión (año 3).} Con la plantilla de diez personas operando a plena capacidad para gestionar el \textit{marketplace}, la nómina asciende a USD 70.000 en sueldos y USD 14.000 en cargas sociales. Esta nómina comprende a los ocho profesionales incorporados y, a partir del punto de equilibrio (alrededor del mes 27), a los socios fundadores.
\end{itemize}
La remuneración de los socios fundadores se difiere hasta alcanzar el punto de equilibrio.

\subsection{Presupuesto financiero}
El presupuesto financiero anual resulta de descontar a los ingresos la totalidad de los egresos y el crecimiento de la plantilla, frente a un aumento de suscriptores de 986 a 2.775 (importes en USD, antes de impuestos):

\begin{center}
\small
\begin{tabularx}{\textwidth}{|X|r|r|r|r|}
\hline
 & \textbf{Año 0} & \textbf{Año 1} & \textbf{Año 2} & \textbf{Año 3}\\ \hline
Ingresos & & 35.000 & 95.000 & 195.000\\ \hline
Costos variables (33\,\%) & & $-$11.550 & $-$31.350 & $-$64.350\\ \hline
Costos fijos operativos & & $-$15.400 & $-$15.900 & $-$16.900\\ \hline
Sueldos & & 0 & $-$24.000 & $-$70.000\\ \hline
Cargas sociales & & $-$2.500 & $-$6.000 & $-$14.000\\ \hline
Inversión inicial & $-$27.610 & & & \\ \hline
\textbf{Flujo neto de caja} & \textbf{$-$27.610} & \textbf{5.550} & \textbf{17.750} & \textbf{29.750}\\ \hline
\textbf{Flujo acumulado} & $-$27.610 & $-$22.060 & $-$4.310 & 25.440\\ \hline
\end{tabularx}
\end{center}

El descuento de los flujos utiliza una tasa de corte del 9\,\%, correspondiente a una sociedad comercial financiada exclusivamente con fondos propios.

\subsection{Evaluación de la inversión}
Bajo los supuestos del escenario base, y antes de impuestos, la inversión de USD 27.610 muestra una rentabilidad alta a tres años:
\begin{itemize}
  \item \textbf{Valor actual neto (VAN):} USD 15.394,02. El resultado positivo indica que el proyecto genera valor económico por encima del costo de oportunidad.
  \item \textbf{Tasa interna de retorno (TIR):} 31,41\,\%, superior a la tasa de corte del 9\,\%, lo que compensa el riesgo de una expansión rápida de la plantilla.
  \item \textbf{Período de repago:} la inversión se recupera alrededor del mes 26, es decir, en los primeros meses del año 3, y el horizonte cierra con un saldo acumulado de USD 25.440.
\end{itemize}

\subsubsection{Punto de equilibrio operativo}
El punto de equilibrio anual es el nivel de ventas que cubre los costos fijos con el margen de contribución del 67\,\% (ingresos menos costos variables). Se resume en la tabla siguiente.

\begin{center}
\small
\begin{tabularx}{\textwidth}{|X|r|r|r|}
\hline
 & \textbf{Año 1} & \textbf{Año 2} & \textbf{Año 3}\\ \hline
Costos fijos anuales (USD) & 17.900 & 45.900 & 100.900\\ \hline
Ventas de equilibrio (USD) & 26.716 & 68.507 & 150.597\\ \hline
Suscriptores de equilibrio & 763 & 1.370 & 2.151\\ \hline
Suscriptores planificados & 986 & 1.889 & 2.775\\ \hline
Margen de seguridad & 23\,\% & 27\,\% & 22\,\%\\ \hline
\end{tabularx}
\end{center}

\subsection{Escenarios de riesgo}
Se evalúan tres escenarios sobre las previsiones iniciales, definidos por una variación en la cantidad de suscriptores (antes de impuestos):
\begin{itemize}
  \item \textbf{Optimista (efecto red inmediato):} los creadores de contenido del \textit{marketplace} viralizan la plataforma en canales de transmisión y reducen el costo de marketing. Se supone un 20\,\% más de suscriptores en todos los años.
  \item \textbf{Pesimista 1 (fricción al precio):} el escalonamiento del precio de la suscripción premium genera resistencia en el público latinoamericano. La plataforma atrae DMs, pero estos permanecen en el plan gratuito, lo que carga la infraestructura sin generar ingresos y consume la reserva de capital de trabajo. Se supone un 30\,\% menos de suscriptores en todos los años.
  \item \textbf{Pesimista 2 (reacción de los competidores):} D\&D Beyond o Roll20 lanzan herramientas gratuitas específicas para partidas presenciales y en español, lo que bloquea la adopción esperada. Se supone un 40\,\% menos de suscriptores desde el año 2.
\end{itemize}

\begin{center}
\small
\begin{tabularx}{\textwidth}{|X|r|r|r|r|r|}
\hline
\textbf{Escenario} & \textbf{Año 1} & \textbf{Año 2} & \textbf{Año 3} & \textbf{VAN (9\,\%)} & \textbf{TIR}\\ \hline
Base & 5.550 & 17.750 & 29.750 & 15.394 & 31,4\,\%\\ \hline
Optimista & 10.240 & 30.480 & 55.880 & 50.589 & 70,9\,\%\\ \hline
Pesimista 1 & $-$1.485 & $-$1.345 & $-$9.445 & $-$37.398 & n/d\\ \hline
Pesimista 2 & 5.550 & $-$7.710 & $-$22.510 & $-$46.389 & n/d\\ \hline
\end{tabularx}
\end{center}
\small Flujos netos anuales en USD. n/d: la TIR no existe porque el proyecto no recupera la inversión.\normalsize

\medskip
Los resultados muestran que el proyecto es sensible a la cantidad de suscriptores: una desviación de apenas el 10\,\% por debajo del escenario base ya vuelve levemente negativo el VAN. En el escenario optimista la inversión se recupera alrededor del mes 19. En los pesimistas el proyecto no recupera la inversión en el horizonte; en el primero, el déficit operativo acumulado (USD 12.275) queda cubierto por la reserva de capital de trabajo, mientras que en el segundo (USD 24.670) la supera levemente.

\subsection{Plan de contingencia}
Ante una desviación negativa del escenario base, como un retraso en el volumen de suscripciones, la empresa:
\begin{itemize}
  \item prolonga de manera indefinida la postergación de la remuneración de los socios fundadores;
  \item reduce a la mitad el presupuesto variable de marketing, limitando el alcance a la publicidad orgánica y a las recomendaciones directas en eventos;
  \item congela la contratación de servicios en la nube destinados al escalado.
\end{itemize}
Solo con la reducción del marketing, el déficit operativo acumulado se reduce de USD 24.670 a USD 8.995 en el escenario pesimista 2 y de USD 12.275 a un superávit de USD 4.786 en el pesimista 1, aunque en ninguno de los dos casos se recupera la inversión inicial dentro del horizonte.

\subsection{Plan de salida}
Ante una contingencia insalvable al cierre del proyecto, se suspenden los débitos automáticos y se emite un aviso de 60 días a los usuarios para que descarguen localmente sus campañas. Luego se liquida la totalidad del equipamiento de red y procesamiento a su valor residual, y el capital resultante se destina a cubrir las obligaciones legales finales y los costos de disolución societaria.
